\documentclass{article}
\usepackage[T1]{fontenc}
\usepackage{lmodern}
\usepackage{graphicx}
\usepackage{amsmath,amssymb}
\usepackage[a4paper,margin=2cm]{geometry}
\usepackage[absolute,overlay]{textpos}
\setlength{\TPHorizModule}{1mm}
\setlength{\TPVertModule}{1mm}
\usepackage[indonesian]{babel}
\usepackage{hyperref}
\setlength{\emergencystretch}{2em}
% unit: Halaman 34

\begin{document}

% === Halaman 34 (python/native) ===

Cipher-Feedback (CFB)

• Mengatasi kekurangan pada mode CBC apabila diterapkan pada 

pengiriman data yang belum mencapai ukuran satu blok  

• Data dienkripsi dalam unit yang lebih kecil (s) daripada ukuran blok (b). 

• Unit data  yang dienkripsi panjangnya bisa 1 bit, 2 bit, 4-bit, 8 bit, dan lain-

lain. 

• Bila unit yang dienkripsi adalah satu karakter setiap kalinya, maka mode 

CFB-nya disebut CFB 8-bit. 

34

\end{document}
